\documentclass[12pt,a4paper,oneside,ngerman]{article}
%\usepackage[utf8]{inputenc}
\usepackage[T1]{fontenc}
\usepackage[default]{raleway}
\usepackage[a4paper, left=25mm, right=25mm, top=30mm, bottom=25mm, headheight=30mm]{geometry}
%\usepackage{csquotes}
\usepackage{ngerman}
\usepackage{textcomp}
\usepackage{comment}
\usepackage{enumerate}
\usepackage{graphicx}
\usepackage{fancyhdr}
\usepackage{epstopdf}
\usepackage{titling}
\usepackage[backend=biber,style=apa]{biblatex}
\DeclareFieldFormat{apacase}{#1}
\DeclareLanguageMapping{ngerman}{ngerman-apa}
\usepackage[singlespacing]{setspace}
\usepackage[plainpages=false,pdfborder={0 0 0}]{hyperref}
\newcommand{\inlinequote}[2]{\glqq{}#1\grqq{} #2}
\graphicspath{{/Users/wombat/nextcloud/ba-geometry/logos/}}

%--------------------------------------------------------------------------------------------------------------
% pagestyles
%--------------------------------------------------------------------------------------------------------------

\fancyhf{}
\fancyhead[L]{Pädagogische Hochschule Heidelberg\\Prof. Dr. Christian Spannagel\\spannagel@ph-heidelberg.de}
\fancyhead[R]{\includegraphics[height=15mm]{phhd.eps}}
\fancyfoot[C]{Seite \thepage}
\fancypagestyle{docstyle}{%
   \fancyhf{}
   \renewcommand{\headrulewidth}{0pt}
   \fancyfoot[C]{Seite \thepage}
}

\setlength{\parskip}{1 em} % fügt Leerzeile zwischen Absätzen ein
\setlength{\parindent}{0 em}


%--------------------------------------------------------------------------------------------------------------
% Basic meta data
%--------------------------------------------------------------------------------------------------------------

\title{Definieren}
\author{Prof. Dr. Christian Spannagel}
\date{\today}
\hypersetup{%
    pdfauthor={\theauthor},%
    pdftitle={\thetitle},%
    pdfsubject={Aufgabenblatt Geometrie},%
    pdfkeywords={geomerie}
}

%\bibliography{/Users/wombat/data/latex/literatur.bib}
%\bibliographystyle{plain}
%\defbibheading{references}[\refname]{\subsection{#1}
%   \markboth{#1}{#1}}


%--------------------------------------------------------------------------------------------------------------
% document
%--------------------------------------------------------------------------------------------------------------

\begin{document}
\pagenumbering{arabic}
\pagestyle{fancy}
\begin{center}
\vspace*{5mm}
{\LARGE\textbf{\thetitle}}
\end{center}

\vspace*{10mm}

\vspace*{10mm}

\textbf{Aufgabe 1 (Vierecke):}  Definiert die folgenden Begriffe, indem ihr die Lücken füllt. Findet, wenn es geht, mehrere Möglichkeiten!
\begin{enumerate}[a)]
\item Ein Rechteck mit\ldots ist ein Quadrat.
\item Ein Parallelogramm\ldots ist ein Rechteck.
\item Ein Trapez\ldots
\item Ein Viereck\ldots
\item Es seien $A$, $B$, $C$ und $D$ vier Punkte, von denen drei nicht auf ein und derselben Geraden liegen. Ein Viereck ist...
\end{enumerate}

\textbf{Aufgabe 2 (Vierecksmengen):} Es sei $V$ die Menge aller (konvexen) Vierecke. Ferner seien: 
\begin{itemize}
\item $T$ die Menge aller Trapeze,
\item $P$ die Menge aller Parallelogramme,
\item $S$ die Menge aller symmetrischen Trapeze,
\item $R$ die Menge aller Rechtecke,
\item $Q$ die Menge aller Quadrate,
\item $R_a$ die Menge aller Rauten und
\item $D$ die Menge aller Drachen.
\end{itemize}
\begin{enumerate}
\item Gebt zwei Vierecksmengen $A$ und $B$ an, für die $A\cup B = A$ gilt.
\item Bestimmt $R \cap Q$.
\item Maxi behauptet: $R_a \cup R = P$ . Stimmt das?
\item Definiert, was man unter einem Element von $S$ versteht.
\item Bestimmt $D\cap T$.
\end{enumerate}

\textbf{Aufgabe 3 (Dreiecke):}  Definiert die folgenden Begriffe: Dreieck, gleichschenkliges Dreieck, gleichseitiges Dreieck, spitzwinkliges Dreieck, rechtwinkliges Dreieck, stumpfwinkliges Dreieck

\textbf{Aufgabe 4 (Kreise):} Definiert die folgenden Begriffe: Kreis, Radius, Durchmesser.

\textbf{Aufgabe 5 (Definieren):}  Reflektiert einmal, was ihr bei euren Definitionen gemacht habt. Wie funktioniert Definieren? Was muss man beachten?

\textbf{Aufgabe 6 (Geogebra):} Installiert euch Geogebra und macht euch mit den geometrischen Konstruktionswerkzeugen vertraut. Konstruiert ein Parallelogramm, ein Quadrat, ein gleichschenkliges und ein gleichseitiges Dreieck jeweils auf verschiedene Arten.


%\newpage
\vspace*{10mm}
\includegraphics[height=10mm]{cc-zero.png}

\tiny{This document is subject to the Creative Commons Zero (CC0) License.
\\ To create this document, we used \LaTeX{}.

Christians Social Media Kanäle:\\
BeReal: \url{https://bere.al/cspannagel}\\
Bluesky: \url{https://bsky.app/profile/cspannagel.bsky.social}\\
Discord: \url{https://tinyurl.com/chrisp-discord}\\
Facebook: \url{https://www.facebook.com/chr.spannagel}\\
Instagram: \url{https://www.instagram.com/dunkelmunkel/}\\
LinkedIn: \url{https://www.linkedin.com/in/christian-spannagel-31b7b4a1/}\\
Mastodon: \url{https://scholar.social/@cspannagel}\\
Snapchat: \url{https://www.snapchat.com/add/spannagelc}\\
Telegram: \url{https://t.me/cspannagel}\\
Tellonym: \url{https://tellonym.me/christian.spannagel}\\
TikTok: \url{https://www.tiktok.com/@_cspannagel_}\\
Threads: \url{https://www.threads.net/@dunkelmunkel}\\
Twitch: \url{https://www.twitch.tv/cspannagel}\\
X/Twitter: \url{https://twitter.com/dunkelmunkel}\\
Youtube: \url{https://www.youtube.com/pharithmetik}}\\


%\pagestyle{docstyle}
\end{document}
